% ── Použité technológie ───────────────────────────────────────
\chapter{Použité technológie}

        Aplikáciu sme postavili na modernom zásobniku technológií, ktorý pokrýva backendové
spracovanie požiadaviek, frontendové zobrazenie a~správu databázy. Výber technológií sme
podmienili požiadavkou na lokálnu prevádzku, jednoduchosť nasadenia a~dostupnosť
kvalitných knižníc pre AI~integráciu.

\section{Backend}

        Backendovú vrstvu aplikácie sme implementovali v~jazyku Python s~využitím
microframeworku Flask. Táto kombinácia umožňuje rýchly vývoj, prehľadnú štruktúru kódu
a~jednoduchú integráciu s~knižnicami pre spracovanie dokumentov a~umelú inteligenciu.

\subsection{Python}

        Python je vysokoúrovňový, interpretovaný programovací jazyk známy svojou
jednoduchosťou a~čitateľnosťou kódu~\cite{python_org}. Bol vytvorený v~roku 1991 Guidom van Rossum a~v~súčasnosti
je jedným z~najpopulárnejších jazykov na svete. Python vyniká minimalistickým syntaxom
a~obrovským ekosystémom knižníc pre rôzne oblasti aplikácií od webových aplikácií cez
dátovú analýzu až po umelú inteligenciu a~strojové učenie.

        V~kontexte webových aplikácií Python ponúka mocné frameworky ako Django a~Flask,
ktoré umožňujú rýchly vývoj robustných backendových riešení. Python je tiež populárny
v~oblasti NLP (Natural Language Processing), čím je ideálny pre náš projekt využívajúci
AI~chatbot.

        \textbf{Výhody Pythonu:}

\begin{itemize}
  \item Jednoduchá a~čitateľná syntax -- Python kód je ľahko zrozumiteľný, čo urýchľuje vývoj
  \item Obrovský ekosystém knižníc -- pre takmer akúkoľvek úlohu existuje kvalitná knižnica
  \item Výkon v~AI/ML -- najobľúbenejší jazyk pre machine learning a~NLP aplikácie
  \item Aktívna komunita -- veľké množstvo tutoriálov, dokumentácie a~podpory
\end{itemize}

        \textbf{Nevýhody Pythonu:}

\begin{itemize}
  \item Pomalší ako kompilované jazyky -- interpretovaný jazyk má nižší výkon ako
        nízkoúrovňové jazyky ako C++ alebo Java
  \item Väčšia spotreba pamäte -- dynamické typovanie a~runtime overhead
  \item GIL (Global Interpreter Lock) -- v~multi-threaded aplikáciách limituje paralelizmus
\end{itemize}

        Pre náš projekt boli výhody Pythonu, najmä jednoduchosť a~vhodnosť pre AI~integráciu,
podstatne dôležitejšie ako nevýhody.

\subsection{Flask Framework}

        Flask je ľahký a~flexibilný Python web framework, ktorý je ideálny pre vývoj
webových aplikácií~\cite{flask_docs}. Na rozdiel od Django, ktorý je \emph{all-in-one} riešenie, ponúka Flask
minimalistický prístup, ktorý umožňuje vývojárom zvoliť si vlastné knižnice a~komponenty.

        Flask je postavený na Werkzeug (WSGI toolkit)~\cite{werkzeug_docs} a~Jinja2 (template engine)~\cite{jinja2_docs} a~poskytuje
základné funkcionality ako:

\begin{itemize}
  \item Routing -- jednoduché definovanie URL ciest
  \item Request/Response handling -- spracovanie bezpečných HTTPS požiadaviek
  \item Session management -- správa používateľských relácií s~HTTPS cookies
  \item Error handling -- ošetrovanie chýb
\end{itemize}

        V~našom projekte sme Flask zvolili pre jeho jednoduchosť, flexibilnosť a~možnosť
ľahkej integrácie s~rôznymi knižnicami pre AI~a~databázové operácie. Aplikácia beží
na HTTPS s~SSL/TLS šifrovaním pre zabezpečenú komunikáciu.

\subsection{Použité knižnice a moduly}

        V~projekte sme využili viaceré Python knižnice:

\begin{itemize}
  \item \textbf{SQLAlchemy}~\cite{sqlalchemy_docs} -- ORM (Object-Relational Mapping) knižnica pre komunikáciu s~databázou
  \item \textbf{Flask-SQLAlchemy}~\cite{sqlalchemy_docs} -- integrácia SQLAlchemy s~Flask
  \item \textbf{Flask-Login}~\cite{flask_login} -- spravovanie používateľských relácií a~autentifikácie
  \item \textbf{Werkzeug}~\cite{werkzeug_docs} -- WSGI toolkit, ktorý Flask používa na spracovanie HTTPS požiadaviek
        s~SSL/TLS podporou
  \item \textbf{Flask-WTF}~\cite{flask_wtf} -- CSRF ochrana vo formulároch
  \item \textbf{Flask-Caching}~\cite{flask_caching} -- in-memory cache pre zrýchlenie opakovaných požiadaviek
  \item \textbf{ollama}~\cite{ollama_github} -- integrácia s~lokálnym LLM serverom pre AI~chat funkcionalitu
  \item \textbf{google-genai}~\cite{google_genai} -- integrácia s~Google Generative AI API pre náročnejšie úlohy
  \item \textbf{python-docx}~\cite{python_docx} -- čítanie a~písanie Word dokumentov (.docx)
  \item \textbf{PyPDF2}~\cite{pypdf2_docs} -- spracovanie PDF súborov
  \item \textbf{python-dotenv}~\cite{python_dotenv} -- správa citlivých informácií (API kľúče, databázové pripojenia)
  \item \textbf{Jinja2}~\cite{jinja2_docs} -- template engine pre dynamické HTML stránky
\end{itemize}

        Tieto knižnice tvoria základ backendovej infraštruktúry a~umožňujú efektívne
spracovávanie požiadaviek, správu dát a~integráciu AI~funkcionalít.

\section{Frontend}

        Frontend aplikácie sme postavili na moderných webových technológiách, ktoré
zabezpečujú interaktívne a~responzívne používateľské rozhranie. Hlavným zámerom
bolo vytvoriť používateľsky intuitívne rozhranie s~podporou tmavého režimu.

\subsection{HTML a Jinja2 Šablóny}

        HTML tvorí kostru aplikácie a~je dynamicky generovaný pomocou Jinja2 template
engine-a. Jinja2 je mocný template engine pre Python Flask, ktorý umožňuje:

\begin{itemize}
  \item Dynamické vkladanie dát z~backendu do HTML stránok
  \item Opakujúce sa bloky kódu -- \texttt{for} cykly a~\texttt{if} výrazy
  \item Dedičnosť šablón -- \texttt{extends} a~\texttt{block} direktívy
\end{itemize}

        Aplikáciu sme navrhli s~modulárnou štruktúrou šablón s~hlavným \texttt{base.html} súborom,
z~ktorého dedia všetky ďalšie šablóny. Komponenty ako header, sidebar, modálne
dialógy a~karty sme oddelili do samostatných šablón pre lepšiu udržiavateľnosť.

\subsection{JavaScript a Interaktivita}

        JavaScript zabezpečuje dynamické správanie aplikácie bez nutnosti obnovy celej
stránky. Kľúčové JavaScript moduly, ktoré sme vytvorili:

\begin{itemize}
  \item \textbf{main.js} -- základná logika: dark mode, sidebar, event listeners
  \item \textbf{chat.js} -- interaktívny AI~chat s~real-time update-mi
  \item \textbf{kanban.js} -- drag-and-drop funkcionalita pre Kanban tabuľu
  \item \textbf{prefetch.js} -- prefetching stránky pre rýchlejšiu navigáciu
\end{itemize}

        \textbf{Implementované funkcionality:}

\begin{itemize}
  \item Dark mode -- prepínanie medzi svetlým a~tmavým režimom s~uložením v~localStorage
  \item Sidebar navigácia -- responzívny sidebar s~overlay na mobilných zariadeniach
  \item Živá notifikácia -- dynamické aktualizácie bez obnovenia strany
  \item Validácia formulárov -- klientska aj serverová validácia
\end{itemize}

\subsection{Responzívny Dizajn}

        Aplikáciu sme navrhli ako úplne responzívnu, optimalizovanú pre rôzne veľkosti obrazovky:

\begin{itemize}
  \item \textbf{Mobilné zariadenia} ({\textless}768\,px) -- sidebar sa skrýva, navigácia je zložená
  \item \textbf{Tablety} (768\,px -- 1024\,px) -- upravené rozloženie s~menej priestorom
  \item \textbf{Desktopy} ({\textgreater}1024\,px) -- plné rozloženie so všetkými funkciami
\end{itemize}

\section{Databáza}

        V~aplikácii sme použili relačnú databázu pre trvalé uloženie všetkých dát. Model databázy
sme definovali pomocou SQLAlchemy ORM, ktorý umožňuje pracovať s~databázovými
tabuľkami ako s~Python objektmi.

\subsection{SQLite}

        Ako primárny databázový systém sme zvolili SQLite~\cite{sqlite_org}. SQLite je ľahký, spoľahlivý
a~vhodný pre aplikácie s~malými až strednými dátovými objemami. Nevyžaduje samostatný
databázový server -- celá databáza je uložená v~súbore \texttt{promat.db}.

        \textbf{Výhody SQLite:}

\begin{itemize}
  \item Bez potreby externého servera -- databáza je jednoduchý súbor
  \item Jednoduchá inštalácia a~konfigurácia
  \item Dostatočný výkon pre väčšinu prípadov
  \item ACID transakcie zabezpečujúce integritu dát
\end{itemize}

\section{Ostatné nástroje}

        Okrem hlavných technológií sme využili niekoľko ďalších nástrojov a~služieb
pre rozšírenie funkcionality.

\subsection{AI a LLM Integrácia}

        Aplikáciu sme integrovali s~viacerými AI~modelmi:

\begin{itemize}
  \item \textbf{Ollama}~\cite{ollama_github} -- lokálny LLM server s~open-source modelmi (napr. Gemma~4)
  \item \textbf{Google Generative AI}~\cite{google_genai} -- cloudové AI~API pre náročnejšie úlohy
\end{itemize}

        Tieto integrácie umožňujú AI~chat funkcionalitu a~automatické generovanie informácií.
Predvolený model je \texttt{gemma4:26b}~\cite{gemma4_google} bežiaci lokálne, čo zabezpečuje plnú súkromnosť
spracovávaných dokumentov.

\subsection{Knižnice pre Dokumenty}

        Pre spracovanie dokumentov sme použili:

\begin{itemize}
  \item \textbf{python-docx}~\cite{python_docx} -- čítanie a~písanie Word dokumentov (.docx)
  \item \textbf{PyPDF2}~\cite{pypdf2_docs} -- spracovanie PDF súborov
\end{itemize}

        Pomocou týchto knižníc aplikácia vie extrahovať text z~dokumentov a~umožniť
AI~modelom pracovať s~ich obsahom.

\subsection{Verzia Kontrola a Deployment}

        Na verzionovanie kódu a~správu vývojového procesu sme použili Git~\cite{git_scm}. Každú
funkčnú zmenu sme zaznamenali v~samostatnom commite s~popisom vykonaných úprav.

\subsection{Bezpečnosť a Overovanie}

\begin{itemize}
  \item \textbf{Flask-Login}~\cite{flask_login} -- session management a~autentifikácia
  \item \textbf{Flask-WTF}~\cite{flask_wtf} -- CSRF ochrana vo formulároch
  \item \textbf{python-dotenv}~\cite{python_dotenv} -- správa citlivých informácií (API keys, databázové pripojenia)
\end{itemize}

\subsection{Výkon a Caching}

\begin{itemize}
  \item \textbf{Flask-Caching}~\cite{flask_caching} -- in-memory cache pre zrýchlenie opakovaných požiadaviek
        s~predvoleným časovým limitom 300~sekúnd
\end{itemize}
